% Generated draft: requires mathematical and source review.
\documentclass[11pt]{article}
\usepackage[margin=25mm]{geometry}
\usepackage{fontspec}
\setmainfont{Latin Modern Roman}
\usepackage{amsmath,amssymb,mathtools,unicode-math}
\setmathfont{Latin Modern Math}
\usepackage{enumitem,needspace,xurl,hyperref}
\setlength{\parindent}{0pt}
\setlength{\parskip}{5pt}
\setlength{\emergencystretch}{3em}
\newcommand{\N}{\mathbb{N}}
\newcommand{\Z}{\mathbb{Z}}
\newcommand{\Q}{\mathbb{Q}}
\newcommand{\R}{\mathbb{R}}
\newcommand{\C}{\mathbb{C}}
\newcommand{\F}{\mathbb{F}}
\newcommand{\A}{\mathbb{A}}
\newcommand{\PP}{\mathbb{P}}
\newcommand{\E}{\mathbb{E}}
\newcommand{\eps}{\varepsilon}
\newcommand{\dd}{\,\mathrm{d}}
\begin{document}
\section*{UnsolvedMath ID: 11100010}
\textbf{Source problem code:} catalogue entry
\subsection{Definitions and mathematical statement}
\textbf{Major problem 10 (Mark Hovey).} Establish the intended permanence of the classes $h_0h_i$ in the Adams spectral sequence at every prime $p>3$, despite the flaw in the proof attributed by the source to R. Cohen--Goerss.

In the customary mod-$p$ sphere interpretation, the classical Adams spectral sequence has
\[
 E_2^{s,t}=\operatorname{Ext}_{\mathcal A_p}^{s,t}(\F_p,\F_p)
 \ \Longrightarrow\ \pi_{t-s}(S^0)^{\wedge}_p,
\]
where $\mathcal A_p$ is the mod-$p$ Steenrod algebra. The symbols $h_i$ are the source's Adams classes, and $h_0h_i$ is their product on the appropriate page. To be a nonzero permanent cycle means surviving all later differentials and incoming boundaries to give a nonzero $E_\infty$ class, not merely being an initial-page cycle.

The uniform target is: for every prime $p>3$ and every intended index $i$ for which the specified nonzero product is defined, $h_0h_i$ survives. The excerpt does not define the $h_i$, give their bidegrees or index range, or explicitly name the spectrum and Adams convention. The sphere spectral sequence above is therefore a customary interpretation requiring confirmation, not an additional supplied hypothesis. The summary's phrase ``all $i$'' cannot repair these omissions or make zero/undefined products into substantive permanence assertions.

Hovey reports that Minami found a flaw, apparently fatal to the cited R. Cohen--Goerss proof, and calls the problem still open at the time of the list. This is historical source testimony, not an independent present-day review of that proof. The summary mentions $\beta$ and $\alpha\beta$ families and partial progress, but supplies neither an exact comparison to $h_0h_i$ nor citations establishing the required uniform statement.
\subsection{Short English statement}
Prove that the intended $h_0h_i$ classes survive the Adams spectral sequence at every prime greater than three. The source reports a fatal proof gap, while leaving the class convention and index range unstated.
\subsection{Sources}
\begin{itemize}
\item Local UnsolvedMath registry, permanent ID 11100010, source code catalogue entry; attributed to Mark Hovey, no proposal year supplied.
\item \emph{Hovey's Algebraic Topology Problem List}, Major problems, item 10; surviving mirror \url{https://www-users.cse.umn.edu/~tlawson/hovey/}. The registry records HTML-ordered-list extraction and access on 2026-07-29; the mirror has 65 top-level items on nine category pages.
\item The source names R. Cohen--Goerss for the flawed proof and Minami for finding the gap, without precise bibliographic citations. The summary names Ravenel and others for related families, also without a specific reference.
\end{itemize}
\end{document}
